\documentclass[11pt, a4paper]{report}

% ========== Packages ==========
\usepackage{fontspec}
\usepackage{kotex}
\usepackage{amsmath, amssymb, amsthm}
\usepackage{mathtools}
\usepackage{bm}
\usepackage{graphicx}
\usepackage{booktabs}
\usepackage{array}
\usepackage{tabularx}
\usepackage{longtable}
\usepackage{hyperref}
\usepackage{cleveref}
\usepackage{geometry}
\usepackage{fancyhdr}
\usepackage{titlesec}
\usepackage{enumitem}
\usepackage{xcolor}
\usepackage{tcolorbox}
\usepackage{algorithm}
\usepackage{algorithmic}
\usepackage{tikz}
\usetikzlibrary{arrows.meta, positioning, calc, decorations.pathreplacing, shapes.geometric, fit, patterns, angles, quotes}
\usepackage{pgfplots}
\pgfplotsset{compat=1.18}
\usepackage{caption}
\usepackage{subcaption}
\usepackage{float}
\usepackage{multirow}
\usepackage{siunitx}

\geometry{margin=2.5cm, top=3cm, bottom=3cm}

% ========== Colors ==========
\definecolor{defblue}{RGB}{30, 80, 150}
\definecolor{thmgreen}{RGB}{20, 110, 60}
\definecolor{noteorange}{RGB}{200, 100, 0}
\definecolor{lightblue}{RGB}{230, 240, 255}
\definecolor{lightgreen}{RGB}{230, 255, 235}
\definecolor{lightorange}{RGB}{255, 245, 230}
\definecolor{lightgray}{RGB}{245, 245, 245}
\definecolor{lightred}{RGB}{255, 235, 235}
\definecolor{darkred}{RGB}{180, 30, 30}
\definecolor{biogreen}{RGB}{10, 130, 80}
\definecolor{cvpurple}{RGB}{100, 40, 150}
\definecolor{injuryred}{RGB}{200, 40, 40}

% ========== Theorem-like environments ==========
\theoremstyle{definition}
\newtheorem{definition}{Definition}[chapter]
\newtheorem{example}{Example}[chapter]
\theoremstyle{remark}
\newtheorem{remark}{Remark}[chapter]

% ========== tcolorbox ==========
\newtcolorbox{keyidea}[1][]{colback=lightblue, colframe=defblue, fonttitle=\bfseries, title={Key Idea}, #1}
\newtcolorbox{importantnote}[1][]{colback=lightorange, colframe=noteorange, fonttitle=\bfseries, title={Important}, #1}
\newtcolorbox{summary}[1][]{colback=lightgreen, colframe=thmgreen, fonttitle=\bfseries, title={Summary}, #1}
\newtcolorbox{clinicalnote}[1][]{colback=lightred, colframe=darkred, fonttitle=\bfseries, title={Clinical Note}, #1}
\newtcolorbox{biomechbox}[1][]{colback=lightgreen, colframe=biogreen, fonttitle=\bfseries, title={Biomechanics}, #1}
\newtcolorbox{cvbox}[1][]{colback={cvpurple!5}, colframe=cvpurple, fonttitle=\bfseries, title={Computer Vision}, #1}
\newtcolorbox{injurybox}[1][]{colback={injuryred!8}, colframe=injuryred, fonttitle=\bfseries, title={Injury Risk}, #1}
\newtcolorbox{setupbox}[1][]{colback=lightgray, colframe=gray!60, fonttitle=\bfseries, title={Setup}, #1}

% ========== Custom commands ==========
\newcommand{\vx}{\mathbf{x}}
\newcommand{\vd}{\mathbf{d}}
\newcommand{\vo}{\mathbf{o}}
\newcommand{\vc}{\mathbf{c}}
\newcommand{\vr}{\mathbf{r}}
\newcommand{\vp}{\mathbf{p}}
\newcommand{\vv}{\mathbf{v}}
\newcommand{\vF}{\mathbf{F}}
\newcommand{\vM}{\mathbf{M}}
\newcommand{\va}{\mathbf{a}}
\newcommand{\vg}{\mathbf{g}}
\newcommand{\vmu}{\bm{\mu}}
\newcommand{\vbeta}{\bm{\beta}}
\newcommand{\vtheta}{\bm{\theta}}
\newcommand{\mSigma}{\bm{\Sigma}}
\newcommand{\mR}{\mathbf{R}}
\newcommand{\mW}{\mathbf{W}}
\newcommand{\mJ}{\mathbf{J}}
\newcommand{\mI}{\mathbf{I}}
\newcommand{\RR}{\mathbb{R}}
\newcommand{\degps}{\,\text{deg/s}}
\newcommand{\Nm}{\,\text{N}\cdot\text{m}}
\newcommand{\BW}{\,\text{BW}}

% ========== Header/Footer ==========
\pagestyle{fancy}
\fancyhf{}
\fancyhead[L]{\leftmark}
\fancyhead[R]{\thepage}
\renewcommand{\headrulewidth}{0.4pt}

% ========== Title formatting ==========
\titleformat{\chapter}[display]
  {\normalfont\huge\bfseries\color{defblue}}
  {\chaptertitlename\ \thechapter}{20pt}{\Huge}
\titleformat{\section}
  {\normalfont\Large\bfseries\color{defblue!80!black}}
  {\thesection}{1em}{}
\titleformat{\subsection}
  {\normalfont\large\bfseries\color{defblue!60!black}}
  {\thesubsection}{1em}{}

\begin{document}

% ==================== TITLE PAGE ====================
\begin{titlepage}
\centering
\vspace*{1.5cm}
{\Large\color{gray} Part 3}
\vspace{0.5cm}

{\Huge\bfseries\color{defblue} 7시점 다시점 촬영 기반\\[0.3cm] 야구 투구 동작 분석}

\vspace{1cm}

{\LARGE\color{defblue!70!black} 240\,Hz 다시점 RGB + Vicon 동기 촬영으로\\3D Gaussian Splatting 기반 투구 생체역학 분석}

\vspace{2cm}

\begin{tikzpicture}[scale=0.65]
% Mound (top view)
\fill[brown!30] (0,0) circle (1.5);
\node[font=\small] at (0,0) {Pitcher};
% Cameras
\foreach \angle/\label in {0/Cam1, 45/Cam2, 90/Cam3, 135/Cam4, 180/{Cam5\,(rear)}, 225/Cam6, 315/Cam7} {
  \fill[defblue] (\angle:5) circle (0.2);
  \node[font=\tiny, text=defblue] at (\angle:6) {\label};
  \draw[-Stealth, gray, thin] (\angle:5) -- (\angle:1.8);
}
% Home plate
\fill[gray!50] (270:4) -- ++(0.3,0.2) -- ++(0.3,0) -- ++(0,-0.4) -- ++(-0.3,0) -- cycle;
\node[font=\tiny] at (270:4.8) {Home};
\end{tikzpicture}

\vspace{2cm}
{\large 컴퓨터 공학 $\times$ 스포츠의학}
\vspace{0.5cm}

{\large 2026}

\vspace{\fill}
{\small\color{gray} 본 교재는 Part 1 (NeRF$\rightarrow$GART), Part 2 (3DGS+동작분석)의 후속편으로,\\
야구 투구라는 구체적 시나리오에서의 다시점 분석 가능성을 탐구합니다.}
\end{titlepage}

\tableofcontents
\newpage

% ==================== PREFACE ====================
\chapter*{서문}
\addcontentsline{toc}{chapter}{서문}

이 교재는 매우 구체적인 시나리오를 다룹니다:

\begin{setupbox}[title={실험 세팅}]
\begin{itemize}[leftmargin=*]
  \item \textbf{피사체}: 야구 투수, 오버헤드(overhand) 투구 동작
  \item \textbf{RGB 카메라}: 7대, 모두 240\,Hz
  \begin{itemize}
    \item 6대: 1920$\times$1080 (가로, landscape)
    \item 1대: 1080$\times$1920 (세로, portrait) --- 피사체 뒤쪽
  \end{itemize}
  \item \textbf{Vicon 모션캡처}: 동시 촬영 (250--500\,Hz)
  \item \textbf{목표}: 투구 생체역학 분석 + 3DGS 기반 3D 재구성
\end{itemize}
\end{setupbox}

이 설정에서 ``3DGS 기반 인체 재구성 + 생체역학 분석''이 어떻게 가능한지, 기존 기술(Vicon, 마커리스 상용 시스템)과 비교하여 어떤 이점이 있는지, 그리고 야구 투구 동작의 생체역학적 의미가 무엇인지를 함께 설명합니다.


% ====================================================================
% CHAPTER 1: 투구 생체역학 기초
% ====================================================================
\chapter{야구 투구 생체역학의 기초}
\label{ch:pitching_biomech}

\section{왜 투구 분석이 중요한가?}

야구 투구는 인체에서 가장 빠른 관절 운동을 포함합니다. 어깨 내회전 속도는 \textbf{7,000--9,000\degps}에 달하며, 이는 약 \textbf{122--157\,rad/s}로 인체에서 측정된 가장 빠른 움직임입니다. 이 극단적인 속도와 반복적인 부하가 만나 부상이 빈번하게 발생합니다.

\begin{injurybox}
2023년 기준 \textbf{MLB 현역 투수의 35.3\%}가 토미 존 수술(UCL 재건술)을 받은 경험이 있으며, 이 비율은 2016년 대비 29\% 증가했습니다. 투구 생체역학 분석은 이러한 부상의 위험 인자를 식별하고, 예방 전략을 수립하는 데 핵심적인 도구입니다.
\end{injurybox}


\section{투구 동작의 6단계}
\label{sec:pitching_phases}

\begin{figure}[H]
\centering
\begin{tikzpicture}[
  phase/.style={rectangle, draw, rounded corners, minimum width=1.8cm, minimum height=1.2cm, align=center, font=\scriptsize},
  event/.style={diamond, draw, aspect=2, minimum width=0.5cm, align=center, font=\tiny, fill=noteorange!30},
  arrow/.style={-{Stealth[length=2mm]}, thick}
]
% Phases
\node[phase, fill=lightblue] (p1) at (0,0) {1. 와인드업\\(Wind-up)};
\node[phase, fill=lightblue] (p2) at (3,0) {2. 스트라이드\\(Stride)};
\node[phase, fill=lightorange] (p3) at (6,0) {3. 후기 코킹\\(Late Cocking)};
\node[phase, fill=injuryred!20] (p4) at (9,0) {4. 가속\\(Acceleration)};
\node[phase, fill=injuryred!20] (p5) at (12,0) {5. 감속\\(Deceleration)};
\node[phase, fill=lightgreen] (p6) at (15,0) {6. 팔로스루\\(Follow-through)};

% Events
\node[event] (e1) at (1.5,1.5) {최대\\무릎 높이};
\node[event] (e2) at (4.5,1.5) {발 접지\\(FC)};
\node[event] (e3) at (7.5,1.5) {최대\\외회전\\(MER)};
\node[event] (e4) at (10.5,1.5) {공 릴리스\\(BR)};
\node[event] (e5) at (13.5,1.5) {최대\\내회전};

% Arrows
\draw[arrow] (p1) -- (p2);
\draw[arrow] (p2) -- (p3);
\draw[arrow] (p3) -- (p4);
\draw[arrow] (p4) -- (p5);
\draw[arrow] (p5) -- (p6);

\draw[gray, dashed] (e1) -- (1.5, 0.6);
\draw[gray, dashed] (e2) -- (4.5, 0.6);
\draw[gray, dashed] (e3) -- (7.5, 0.6);
\draw[gray, dashed] (e4) -- (10.5, 0.6);
\draw[gray, dashed] (e5) -- (13.5, 0.6);

% Time annotation
\draw[decorate, decoration={brace, amplitude=5pt, mirror}, thick, darkred] (9,-0.8) -- (12,-0.8)
  node[midway, below=8pt, font=\scriptsize, text=darkred] {\textbf{$\sim$30\,ms (가장 빠른 구간)}};
\end{tikzpicture}
\caption{투구 동작의 6단계와 주요 이벤트. 가속 단계는 약 30\,ms로 전체 동작에서 가장 짧지만, 부하가 가장 큰 구간입니다.}
\label{fig:pitching_phases}
\end{figure}

\subsection{1단계: 와인드업 (Wind-up)}

투수가 앞발을 들어올리는 순간부터 최대 무릎 높이까지. 리듬을 잡고 위치 에너지를 축적하는 단계로, 관절 부하가 상대적으로 낮습니다.

\subsection{2단계: 스트라이드 / 초기 코킹 (Stride / Early Cocking)}

최대 무릎 높이에서 앞발 접지(Foot Contact, FC)까지. 투수가 홈 방향으로 스트라이드하며, 투구 팔은 외전(abduction)과 외회전(external rotation)을 시작합니다.

\begin{biomechbox}[title={핵심 파라미터: 스트라이드}]
\begin{itemize}[leftmargin=*]
  \item \textbf{스트라이드 길이}: 투수 신장의 약 85\% (범위: 74--93\%)
  \item 숙련도가 높을수록 스트라이드가 길어짐 $\rightarrow$ 공 릴리스 지점이 홈에 가까워짐
  \item 스트라이드 방향은 홈 방향에서 $\pm 5^\circ$ 이내가 이상적
\end{itemize}
\end{biomechbox}

\subsection{3단계: 후기 코킹 (Late Cocking)}

앞발 접지에서 \textbf{최대 외회전(Maximum External Rotation, MER)}까지.

\begin{definition}[골반-몸통 분리 (Hip-Shoulder Separation)]
앞발 접지 시점에서 골반 회전각과 몸통(어깨) 회전각의 차이:
\begin{equation}
\theta_{\text{sep}} = \theta_{\text{pelvis}} - \theta_{\text{trunk}}
\end{equation}
골반은 이미 홈 방향으로 $\sim 35^\circ$ 회전했지만 어깨는 아직 닫혀 있어, 평균 $50 \pm 12^\circ$의 분리가 발생합니다. 이 ``탄성 에너지 저장''이 투구 속도의 약 25\%를 설명합니다.
\end{definition}

후기 코킹 말기에 어깨 외회전은 \textbf{170--180$^\circ$}에 도달하며, 이 시점에서:

\begin{center}
\renewcommand{\arraystretch}{1.3}
\begin{tabular}{lcc}
\toprule
\textbf{파라미터} & \textbf{값} & \textbf{임상적 의미} \\
\midrule
어깨 내회전 토크 피크 & $67\Nm$ & 회전근개 부하 \\
\textcolor{darkred}{\textbf{팔꿈치 외반 토크 피크}} & \textcolor{darkred}{\textbf{$64\Nm$}} & \textcolor{darkred}{\textbf{UCL 부하 (핵심 부상 지표)}} \\
어깨 압축력 & $550$--$770$\,N & 관절와순 부하 \\
\bottomrule
\end{tabular}
\end{center}

\subsection{4단계: 가속 (Acceleration)}

MER에서 \textbf{공 릴리스(Ball Release, BR)}까지. 전체 투구에서 \textbf{가장 짧은 구간($\sim$30\,ms)}이지만, 운동 에너지 전달이 집중됩니다.

\begin{keyidea}
가속 단계에서 어깨는 $\sim$80$^\circ$의 내회전을 약 30\,ms 만에 수행합니다. 이때의 각속도 \textbf{7,000--9,000\degps}는 인체에서 측정된 가장 빠른 관절 운동입니다.
\begin{equation}
\omega_{\text{IR}} = \frac{d\theta_{\text{IR}}}{dt} \approx 7{,}000\text{--}9{,}000\degps \approx 122\text{--}157\,\text{rad/s}
\end{equation}
\end{keyidea}

\subsection{5단계: 감속 (Deceleration)}

공 릴리스에서 최대 내회전까지. 투구에서 \textbf{역학적으로 가장 가혹한 구간}입니다.

\begin{center}
\renewcommand{\arraystretch}{1.3}
\begin{tabular}{lc}
\toprule
\textbf{감속 단계 후방 어깨 부하} & \textbf{크기} \\
\midrule
후방 전단력 (posterior shear) & $400$\,N \\
하방 전단력 (inferior shear) & $300$\,N \\
\textbf{압축력 (compressive)} & \textbf{$>1{,}000$\,N} \\
어깨 견인력 (distraction) & 체중의 90--150\% \\
수평 외전 토크 & $97\Nm$ \\
\bottomrule
\end{tabular}
\end{center}

이 엄청난 힘들은 후방 어깨 근육(후면 삼각근, 극하근, 소원근)이 편심성(eccentric) 수축으로 흡수해야 합니다.

\subsection{6단계: 팔로스루 (Follow-through)}

감속이 지속되며 투수가 필딩 자세로 전환합니다.


\section{핵심 부상 메커니즘: UCL 부하 패러독스}
\label{sec:ucl_paradox}

\begin{injurybox}[title={UCL 부하 패러독스 (Loading Paradox)}]
\textbf{문제}: 사체(cadaver) 연구에서 UCL(척골측부인대)의 파단 강도(failure load)는 약 $30$--$34\Nm$. 그런데 실제 투구에서 팔꿈치 외반 토크는 $40$--$120\Nm$에 달합니다.

\begin{equation}
\underbrace{\tau_{\text{UCL failure}} \approx 30\text{--}34\Nm}_{\text{인대 파단 강도}} \quad \ll \quad \underbrace{\tau_{\text{valgus peak}} \approx 64\Nm}_{\text{실제 투구 부하}}
\end{equation}

\textbf{해결}: \textbf{굴근-회내근 복합체(Flexor-Pronator Mass, FPM)}와 이두/삼두의 동시 수축이 내반(varus) 반대 토크를 생성하여 UCL을 보호합니다. 그러나 매 투구마다 UCL에 가해지는 부하는 ``거의 파열에 충분한'' 수준이며, 이것이 \textbf{누적 피로 부상(cumulative fatigue injury)}으로 이어집니다.
\end{injurybox}

이 패러독스가 바로 투구 생체역학 분석의 핵심 동기입니다: 팔꿈치 외반 토크를 정확하게 측정하고, 부하 누적을 관리하면 UCL 부상을 예방할 수 있습니다.


\section{공 릴리스와 구종별 특성}
\label{sec:ball_release}

\begin{center}
\renewcommand{\arraystretch}{1.3}
\begin{tabular}{lccc}
\toprule
\textbf{구종} & \textbf{속도 (mph)} & \textbf{회전수 (RPM)} & \textbf{회전 효율} \\
\midrule
포심 패스트볼 & 94--95 & $\sim$2,300 & 95--100\% \\
투심 패스트볼 & $\sim$93 & $\sim$2,150 & -- \\
커터 & 88--89 & $\sim$2,350 & -- \\
슬라이더 & 84--85 & $\sim$2,450 & 20--35\% \\
커브볼 & $\sim$80 & 2,430--2,530 & $\sim$78\% \\
체인지업 & $\sim$86 & $\sim$1,750 & -- \\
\bottomrule
\end{tabular}
\end{center}

\begin{definition}[회전 효율 (Spin Efficiency)]
총 회전 중 직선 궤적에 기여하는 ``유효 회전''의 비율:
\begin{equation}
\text{Spin Efficiency} = \frac{\omega_{\text{transverse}}}{\omega_{\text{total}}} \times 100\%
\end{equation}
패스트볼은 백스핀이 대부분 유효하므로 95--100\%, 슬라이더는 자이로 회전(gyro spin) 성분이 커서 20--35\%입니다.
\end{definition}


% ====================================================================
% CHAPTER 2: 우리의 실험 세팅
% ====================================================================
\chapter{실험 세팅: 7대 카메라 + Vicon 동시 촬영}
\label{ch:setup}

\section{카메라 배치 설계}
\label{sec:camera_placement}

\begin{figure}[H]
\centering
\begin{tikzpicture}[scale=0.8]
% Mound
\fill[brown!20] (0,0) circle (1.8);
\draw[brown!50] (0,0) circle (1.8);
\node[font=\small\bfseries] at (0,0.3) {투수};
\node[font=\tiny, text=gray] at (0,-0.2) {(투구판)};

% Home plate direction
\draw[-Stealth, thick, gray!50] (0,-2) -- (0,-5) node[right, font=\small] {홈 방향};

% Camera positions
\def\camrad{6}

% Cam 1: Front-left (landscape)
\fill[defblue] (-3,-4.5) circle (0.25);
\node[font=\scriptsize\bfseries, text=defblue] at (-3,-5.2) {Cam 1};
\node[font=\tiny, text=gray] at (-3,-5.6) {1920$\times$1080};
\draw[-Stealth, defblue!50, thin] (-3,-4.5) -- (-0.8,-0.5);

% Cam 2: Front-right (landscape)
\fill[defblue] (3,-4.5) circle (0.25);
\node[font=\scriptsize\bfseries, text=defblue] at (3,-5.2) {Cam 2};
\node[font=\tiny, text=gray] at (3,-5.6) {1920$\times$1080};
\draw[-Stealth, defblue!50, thin] (3,-4.5) -- (0.8,-0.5);

% Cam 3: Left side (landscape)
\fill[defblue] (-5.5,-1) circle (0.25);
\node[font=\scriptsize\bfseries, text=defblue] at (-6.5,-1) {Cam 3};
\node[font=\tiny, text=gray] at (-6.5,-1.4) {1920$\times$1080};
\draw[-Stealth, defblue!50, thin] (-5.5,-1) -- (-1.5,0);

% Cam 4: Right side (landscape)
\fill[defblue] (5.5,-1) circle (0.25);
\node[font=\scriptsize\bfseries, text=defblue] at (6.5,-1) {Cam 4};
\node[font=\tiny, text=gray] at (6.5,-1.4) {1920$\times$1080};
\draw[-Stealth, defblue!50, thin] (5.5,-1) -- (1.5,0);

% Cam 5: Rear (portrait!)
\fill[cvpurple] (0,5.5) circle (0.25);
\node[font=\scriptsize\bfseries, text=cvpurple] at (0,6.3) {Cam 5 (후면)};
\node[font=\tiny, text=cvpurple] at (0,5.9) {\textbf{1080$\times$1920}};
\draw[-Stealth, cvpurple!50, thin] (0,5.5) -- (0,1.5);

% Cam 6: Back-left (landscape)
\fill[defblue] (-4,4) circle (0.25);
\node[font=\scriptsize\bfseries, text=defblue] at (-5.2,4) {Cam 6};
\node[font=\tiny, text=gray] at (-5.2,3.6) {1920$\times$1080};
\draw[-Stealth, defblue!50, thin] (-4,4) -- (-0.8,0.5);

% Cam 7: Back-right / elevated (landscape)
\fill[defblue] (4,4) circle (0.25);
\node[font=\scriptsize\bfseries, text=defblue] at (5.2,4) {Cam 7};
\node[font=\tiny, text=gray] at (5.2,3.6) {1920$\times$1080};
\draw[-Stealth, defblue!50, thin] (4,4) -- (0.8,0.5);

% Vicon cameras (small dots)
\foreach \a in {20,55,90,125,160,200,235,270,305,340} {
  \fill[thmgreen!70] (\a:7) circle (0.12);
}
\node[font=\tiny, text=thmgreen] at (7.5,3) {Vicon IR};
\node[font=\tiny, text=thmgreen] at (7.5,2.5) {(10--12대)};

% Legend
\draw[thick] (-7,-7) rectangle (7.5,-8.5);
\fill[defblue] (-6.5,-7.5) circle (0.15);
\node[right, font=\scriptsize] at (-6.2,-7.5) {RGB 가로 (1920$\times$1080)};
\fill[cvpurple] (-6.5,-8) circle (0.15);
\node[right, font=\scriptsize] at (-6.2,-8) {RGB 세로 (1080$\times$1920)};
\fill[thmgreen!70] (-1.5,-7.5) circle (0.1);
\node[right, font=\scriptsize] at (-1.2,-7.5) {Vicon IR 카메라};
\node[font=\scriptsize\bfseries] at (4.5,-7.7) {모두 240\,Hz};
\end{tikzpicture}
\caption{7대 RGB 카메라 + Vicon 배치 상면도. 세로(portrait) 카메라(Cam 5)는 투수 뒤에서 투구 팔의 전체 수직 궤적을 포착합니다.}
\label{fig:camera_layout}
\end{figure}

\subsection{각 카메라의 역할}

\begin{center}
\renewcommand{\arraystretch}{1.3}
\small
\begin{tabular}{clp{7cm}}
\toprule
\textbf{카메라} & \textbf{위치} & \textbf{포착 대상} \\
\midrule
Cam 1--2 & 전방 좌/우 45$^\circ$ & 스트라이드 착지, 몸통 회전, 공 릴리스, 앞발 각도 \\
Cam 3--4 & 측면 좌/우 & 팔 슬롯, 어깨 외회전/내회전, 무릎 굴곡 (시상면) \\
\textbf{Cam 5} & \textbf{후면 세로} & \textbf{투구 팔의 전체 수직 궤적, 어깨 외전각, arm path} \\
Cam 6--7 & 후방 좌/우 & 팔로스루, 몸통 후면, 뒷발 이탈 \\
\bottomrule
\end{tabular}
\end{center}

\begin{importantnote}[title={왜 후면 카메라가 세로(portrait)인가?}]
투수의 투구 팔은 수직으로 큰 범위를 움직입니다---와인드업에서 머리 위, MER에서 등 뒤, 팔로스루에서 무릎 아래까지. 후면에서 보면 이 수직 동선이 가장 잘 드러나며, \textbf{세로 촬영(1080$\times$1920)}은 이 수직 범위를 최대한 포착합니다. 가로 촬영 시 불필요한 좌우 공간에 해상도가 낭비됩니다.
\end{importantnote}


\section{240\,Hz의 의미: 나이퀴스트 정리와 시간 분해능}
\label{sec:temporal_resolution}

\begin{definition}[나이퀴스트-섀넌 표본화 정리]
주파수 $f_{\max}$까지의 성분을 포함하는 신호를 정확히 복원하려면, 최소 $2f_{\max}$ 이상의 샘플링 레이트가 필요합니다:
\begin{equation}
f_s \geq 2 f_{\max}
\end{equation}
\end{definition}

투구 가속 단계의 운동은 100\,Hz 이상의 주파수 성분을 포함합니다. 따라서 최소 200\,Hz 이상의 촬영이 필요하며, \textbf{240\,Hz는 이 조건을 충족}합니다.

\subsection{프레임 레이트별 시간 분해능 비교}

\begin{center}
\renewcommand{\arraystretch}{1.3}
\begin{tabular}{lccc}
\toprule
\textbf{프레임 레이트} & \textbf{프레임 간격} & \textbf{가속 단계 프레임 수} & \textbf{피크 시 프레임당 회전각} \\
\midrule
30\,fps (일반 영상) & 33.3\,ms & $\sim$1 & $\sim$300$^\circ$ \\
120\,fps (슬로모션) & 8.33\,ms & $\sim$3--4 & $\sim$75$^\circ$ \\
\textbf{240\,fps (본 세팅)} & \textbf{4.17\,ms} & \textbf{$\sim$7} & \textbf{$\sim$37.5$^\circ$} \\
500\,fps (Vicon) & 2.0\,ms & $\sim$15 & $\sim$18$^\circ$ \\
\bottomrule
\end{tabular}
\end{center}

\begin{equation}
\Delta\theta = \omega \cdot \Delta t = 9{,}000\degps \times \frac{1}{240}\,\text{s} = 37.5^\circ \text{ per frame (피크 시)}
\end{equation}

30\,fps에서는 가속 단계 전체가 1 프레임 안에 들어가 분석이 불가능합니다. 240\,fps에서는 7 프레임으로 분석이 가능하며, 이는 MER, 공 릴리스 등 핵심 이벤트를 식별하기에 충분합니다.


\section{Vicon과의 동기화}
\label{sec:synchronization}

\subsection{시간 정렬 문제}

\begin{center}
\renewcommand{\arraystretch}{1.3}
\begin{tabular}{lcc}
\toprule
\textbf{시스템} & \textbf{주파수} & \textbf{프레임 간격} \\
\midrule
RGB 카메라 & 240\,Hz & 4.17\,ms \\
Vicon (옵션 1) & 250\,Hz & 4.0\,ms \\
Vicon (옵션 2) & 500\,Hz & 2.0\,ms \\
\bottomrule
\end{tabular}
\end{center}

240\,Hz와 250\,Hz는 정수비가 아닙니다 (24:25). 정확한 정렬 주기:
\begin{equation}
\text{LCM}(240, 250) = 6{,}000 \quad \Rightarrow \quad \frac{1}{6000}\,\text{s} = 0.167\,\text{ms마다 정확히 정렬}
\end{equation}

이는 비현실적이므로, 실용적 접근:
\begin{itemize}[leftmargin=*]
  \item \textbf{매 0.1초}마다 정확히 정렬 (25 RGB 프레임 = 24 Vicon 프레임 at 250\,Hz)
  \item 그 사이: Vicon 데이터를 RGB 타임스탬프에 \textbf{큐빅 스플라인 보간}
  \item 최대 시간 오차: $\pm 2$\,ms (250\,Hz 기준) 또는 $\pm 1$\,ms (500\,Hz 기준)
\end{itemize}

\subsection{동기화 방법}

\begin{enumerate}[leftmargin=*]
  \item \textbf{하드웨어 동기화 (Genlock)}: Vicon Lock 장치가 마스터 클럭 신호를 생성하여 모든 카메라의 프레임 캡처를 동기화. GPO(General Purpose Output) 포트로 외부 RGB 카메라 트리거.
  \item \textbf{동기 신호 (TTL Trigger)}: 촬영 시작 시 모든 카메라에 동시 트리거 펄스 전송.
  \item \textbf{사후 정렬}: LED 플래시 또는 클래퍼 보드로 시각/청각 동기 신호 기록 후 교차상관(cross-correlation)으로 정렬.
\end{enumerate}

\begin{remark}
Vicon이 더 높은 주파수(250/500\,Hz)로 촬영하므로, RGB 프레임마다 시간적으로 가장 가까운 Vicon 프레임이 존재합니다. 보간보다 ``최근접 Vicon 프레임 사용''도 실용적입니다 ($\pm 2$\,ms 이내).
\end{remark}


\section{공간 보정 (Spatial Calibration)}
\label{sec:calibration}

\subsection{Vicon과 RGB 좌표계 통합}

Vicon과 RGB 카메라는 서로 다른 좌표계를 사용합니다. 이를 통합하려면:

\begin{enumerate}[leftmargin=*]
  \item \textbf{공유 보정 객체}: Vicon 마커가 부착된 체커보드를 양 시스템에 동시에 보이게 배치
  \item \textbf{대응점 획득}: 체커보드 코너의 3D 좌표를 Vicon과 RGB 삼각측량에서 각각 측정
  \item \textbf{강체 변환 계산}: 최소 3개 비공선(non-collinear) 대응점으로 회전 $\mR$과 이동 $\mathbf{t}$ 추정:
  \begin{equation}
    \vx_{\text{Vicon}} = \mR \cdot \vx_{\text{RGB}} + \mathbf{t}
  \end{equation}
  SVD 기반 최소자승 정합:
  \begin{equation}
    \mR^*, \mathbf{t}^* = \arg\min_{\mR, \mathbf{t}} \sum_{i=1}^{N} \| \vx_{\text{Vicon},i} - \mR \vx_{\text{RGB},i} - \mathbf{t} \|^2
  \end{equation}
  \item \textbf{검증}: Vicon 마커 3D 좌표를 RGB 이미지에 재투영, 재투영 오차 $<2$ 픽셀 목표
\end{enumerate}

\subsection{RGB 카메라 개별 보정}

각 RGB 카메라의 \textbf{내부 파라미터(intrinsic)} 보정:
\begin{equation}
\mathbf{K} = \begin{bmatrix} f_x & 0 & c_x \\ 0 & f_y & c_y \\ 0 & 0 & 1 \end{bmatrix}
\end{equation}
+ 렌즈 왜곡 계수 ($k_1, k_2, p_1, p_2, k_3$). 체커보드 다시점 촬영 + Zhang's 방법으로 추정합니다.

\textbf{외부 파라미터(extrinsic)}: 카메라 간 상대 위치/회전. 공유 보정 객체 또는 Vicon 좌표계를 기준으로 추정합니다.

\begin{importantnote}[title={세로 카메라의 보정}]
Cam 5 (1080$\times$1920)는 별도의 내부 파라미터 보정이 필요합니다. 해상도와 화각이 다르기 때문입니다. 그러나 외부 파라미터는 동일한 보정 절차로 추정되며, \textbf{3DGS는 이질적인 카메라 해상도를 본질적으로 지원}합니다---각 카메라가 자체 투영 모델을 가집니다.
\end{importantnote}


\section{데이터 저장량 산정}
\label{sec:storage}

\subsection{비압축 RAW}

\begin{align}
\text{프레임당} &= 1920 \times 1080 \times 3\,\text{bytes} = 6{,}220{,}800\,\text{bytes} \approx 5.93\,\text{MB} \\
\text{카메라당 초당} &= 5.93 \times 240 = 1{,}423\,\text{MB/s} \approx 1.42\,\text{GB/s} \\
\text{7대 합계 초당} &= 9.96\,\text{GB/s} \approx 79.7\,\text{Gbps}
\end{align}

\subsection{압축 포맷별 비교}

\begin{center}
\renewcommand{\arraystretch}{1.3}
\begin{tabular}{lcccc}
\toprule
\textbf{포맷} & \textbf{투구 1회 (2초)} & \textbf{세션 (100투)} & \textbf{10명 연구} & \textbf{품질} \\
\midrule
RAW & $\sim$20\,GB & $\sim$2\,TB & $\sim$20\,TB & 무손실 \\
ProRes 422 HQ & $\sim$2--3\,GB & 200--300\,GB & 2--3\,TB & 거의 무손실 \\
H.264 (고품질) & 100--175\,MB & 10--17.5\,GB & 100--175\,GB & 약간의 손실 \\
H.265 (고품질) & 50--90\,MB & 5--9\,GB & 50--90\,GB & 약간의 손실 \\
\bottomrule
\end{tabular}
\end{center}

\begin{remark}
3DGS 학습에서는 \textbf{광도 손실(photometric loss)}이 목적함수이므로, 압축 아티팩트가 노이즈로 작용합니다. 분석용 아카이브는 ProRes 이상, 3DGS 학습용은 고품질 H.265가 실용적인 절충안입니다.
\end{remark}


% ====================================================================
% CHAPTER 3: Vicon 기반 투구 분석
% ====================================================================
\chapter{Vicon 마커 기반 투구 분석}
\label{ch:vicon_pitching}

\section{투구 분석용 마커 세트}

투구 분석은 상지(upper extremity)에 집중하므로, 전신 + 상지 강화 마커 세트(38--59개)를 사용합니다.

\begin{biomechbox}[title={투구 분석 마커 세트 (상지 강화)}]
\begin{itemize}[leftmargin=*]
  \item \textbf{머리}: 밴드에 4개 마커
  \item \textbf{몸통}: C7, T10, 흉골절흔(IJ), 검상돌기(PX), 견갑골
  \item \textbf{상완}: 견봉(acromion), 중간 상완 추적 클러스터, 내측/외측 상과(epicondyle)
  \item \textbf{전완}: 요골/척골 경상돌기(styloid), 전완 추적 클러스터
  \item \textbf{손}: 제2/5 중수골두, 제3 중수골 기저부
  \item \textbf{골반}: 양측 ASIS, 양측 PSIS
  \item \textbf{하지}: 대퇴 클러스터, 슬관절 내/외측, 경골 클러스터, 족관절 내/외측, 종골, 발가락
  \item \textbf{공}: 2개 마커 (공 속도 및 릴리스 시점 계산)
\end{itemize}
\end{biomechbox}

\section{ISB 관절 좌표계 (어깨/팔꿈치)}

Wu 등 (2005)의 ISB 권고안:

\subsection{견관절 (Glenohumeral Joint)}

\textbf{Y-X-Y 오일러 분해}:
\begin{itemize}[leftmargin=*]
  \item 제1 회전 (Y): \textbf{거상 평면(plane of elevation)} --- 팔이 어느 방향으로 들렸는지 (전방/측방)
  \item 제2 회전 (X): \textbf{거상각(elevation angle)} --- 팔이 얼마나 들렸는지
  \item 제3 회전 (Y): \textbf{축 회전(axial rotation)} --- 외회전/내회전
\end{itemize}

\subsection{주관절 (Elbow Joint)}

\begin{itemize}[leftmargin=*]
  \item 굴곡/신전 (flexion/extension)
  \item 회내/회외 (pronation/supination)
  \item 운반각 (carrying angle)
\end{itemize}


\section{Vicon 투구 분석의 고유한 문제점}
\label{sec:vicon_challenges}

\begin{enumerate}[leftmargin=*]
  \item \textbf{마커 가려짐(Occlusion)}: 와인드업 시 글러브 팔과 몸통이 투구 팔 마커를 가림. 팔로스루 시 몸통 회전으로 동일 현상. 갭 필링(gap filling) 알고리즘(스플라인 보간, 칼만 필터)으로 보완하지만, \textbf{가장 빠른 동작 구간에서 가려짐이 발생}하므로 보간 오차가 큼.

  \item \textbf{고속 각운동}: 500\,Hz에서도 프레임 간 $\sim$18$^\circ$ 회전. 마커가 점이 아닌 줄무늬(streak)로 나타나 중심 추정 정확도 저하.

  \item \textbf{마커 탈락}: 극단적 가속도로 마커가 이동/탈락. 전완과 손 부분에서 특히 빈번.

  \item \textbf{연조직 아티팩트}: 삼두근, 삼각근 위의 마커가 근수축 중 실제 뼈 움직임과 10--20\,mm 이상 차이.

  \item \textbf{넓은 포착 볼륨}: 스트라이드($\sim$1.5--1.7\,m) + 와인드업 회전까지 포함하면 상당히 넓은 공간이 필요.
\end{enumerate}


% ====================================================================
% CHAPTER 4: 마커리스 시스템 현황
% ====================================================================
\chapter{마커리스 투구 분석의 현재}
\label{ch:markerless_pitching}

\section{KinaTrax / Hawk-Eye (MLB 공식 시스템)}

\begin{cvbox}[title={KinaTrax 시스템}]
\begin{itemize}[leftmargin=*]
  \item \textbf{카메라}: 8대, 300\,fps, 경기장 관중석에 설치
  \item \textbf{기술}: 컴퓨터 비전 + 딥러닝 기반 마커리스 모션캡처; 실루엣 기반 인체 모델 피팅
  \item \textbf{배포}: 다수 MLB 구장에 설치; 경기 중 실시간 생체역학 데이터 수집
  \item \textbf{검증}: 2023시즌 5개 팀 51명 투수 데이터 수집; 2025년 다수의 검증 연구 발표
\end{itemize}
\end{cvbox}

\section{마커리스 vs 마커 기반: 투구 분석 정확도}

2025년 검증 연구 (18명 대학 투수, 최대 노력 패스트볼):

\begin{center}
\renewcommand{\arraystretch}{1.3}
\small
\begin{tabular}{p{4cm}ccc}
\toprule
\textbf{지표} & \textbf{Vicon (기준)} & \textbf{Theia3D} & \textbf{KinaTrax} \\
\midrule
관절 위치 오차 (MPJPE) & 기준 ($<$1\,mm) & 52.0$\pm$12.3\,mm & 56.6$\pm$9.4\,mm \\
팔꿈치 위치 오차 & 기준 & $\sim$70--75\,mm & $\sim$70--75\,mm \\
관절 각도 RMS 오차 & 기준 & 2--16$^\circ$ & 2--16$^\circ$ \\
회전 속도 오차 & 기준 & 50--400+\degps & 50--400+\degps \\
스트라이드 길이 & 기준 & CCC $>$ 0.85 & CCC $>$ 0.85 \\
어깨 회전 변수 & 기준 & 큰 변동성 & 큰 변동성 \\
\bottomrule
\end{tabular}
\end{center}

\begin{clinicalnote}
마커리스 시스템의 \textbf{관절 위치 오차 52--57\,mm}는 보행 분석에서는 수용 가능하지만, 투구 분석에서는 중요한 한계입니다. 팔꿈치 외반 토크 계산은 모멘트 팔(moment arm)에 매우 민감하며, 수 mm의 관절 중심 오차가 수 Nm의 토크 오차로 증폭됩니다.

특히 \textbf{어깨 회전 변수}에서 가장 큰 불일치가 관찰됩니다---정확히 투구 분석에서 가장 중요한 변수(MER, 내회전 속도)에서 마커리스의 한계가 드러납니다.
\end{clinicalnote}


\section{OpenPose/MediaPipe의 한계}

\begin{itemize}[leftmargin=*]
  \item 30\,fps 영상용으로 설계 $\rightarrow$ 240\,fps 고속 영상에 직접 적용 시 중복 연산
  \item 모션 블러에 취약: 고속 투구 시 팔이 흐릿한 줄무늬로 나타남
  \item \textbf{7,000\degps 이상의 각속도에서 검증된 바 없음}
  \item 단안 영상 기반으로 3D 깊이 추정에 본질적 한계
  \item 관절별 정확도가 균일하지 않음 (족관절 등 작은 움직임의 관절에서 특히 부정확)
\end{itemize}


% ====================================================================
% CHAPTER 5: 3DGS 기반 투구 분석 파이프라인
% ====================================================================
\chapter{3DGS 기반 투구 분석 파이프라인}
\label{ch:3dgs_pitching}

\section{제안 파이프라인 개요}

\begin{figure}[H]
\centering
\begin{tikzpicture}[
  block/.style={rectangle, draw, rounded corners, minimum width=3cm, minimum height=0.9cm, align=center, font=\small},
  data/.style={rectangle, draw, dashed, rounded corners, minimum width=2.5cm, minimum height=0.6cm, align=center, font=\scriptsize},
  arrow/.style={-{Stealth[length=2.5mm]}, thick},
  node distance=0.8cm
]
% Input row
\node[block, fill=lightorange] (rgb) at (0,0) {7대 RGB\\(240\,Hz)};
\node[block, fill=lightgreen] (vicon) at (6,0) {Vicon\\(250--500\,Hz)};

% Sync
\node[block, fill=lightgray] (sync) at (3,-1.5) {시간/공간 동기화\\및 보정};

% Processing
\node[block, fill=lightblue] (pose2d) at (-2,-3.5) {다시점 2D\\자세 추정};
\node[block, fill=lightblue] (smpl) at (3,-3.5) {다시점 SMPL-X\\피팅};
\node[block, fill=lightgreen] (vicon_ik) at (8,-3.5) {Vicon 마커 기반\\IK (골드 스탠다드)};

% 3DGS
\node[block, fill=defblue!20] (3dgs) at (0.5,-5.5) {3DGS 인체 재구성\\(GART/ExAvatar)};

% Outputs
\node[block, fill=cvpurple!15] (pose_ref) at (-3,-7.5) {정제된 SMPL\\자세 $\vtheta^*$};
\node[block, fill=cvpurple!15] (render) at (4,-7.5) {포토리얼리스틱\\임의 시점 렌더링};

% Biomech
\node[block, fill=biogreen!20] (opensim) at (-3,-9.5) {OpenSim\\IK / ID};
\node[block, fill=biogreen!20] (virtual) at (1,-9.5) {가상 마커\\추출};

% Clinical
\node[block, fill=injuryred!15] (clinical) at (-1,-11.5) {투구 생체역학 지표\\(MER, 외반토크, 팔속도)};

% Validation
\node[block, fill=thmgreen!20] (validate) at (6,-9.5) {Vicon 결과와\\비교 검증};

% Arrows
\draw[arrow] (rgb) -- (sync);
\draw[arrow] (vicon) -- (sync);
\draw[arrow] (sync) -- (pose2d);
\draw[arrow] (sync) -- (smpl);
\draw[arrow] (sync) -- (vicon_ik);
\draw[arrow] (pose2d) -- (smpl);
\draw[arrow] (smpl) -- (3dgs);
\draw[arrow] (rgb) -- (3dgs);
\draw[arrow] (3dgs) -- (pose_ref);
\draw[arrow] (3dgs) -- (render);
\draw[arrow] (pose_ref) -- (opensim);
\draw[arrow] (pose_ref) -- (virtual);
\draw[arrow] (virtual) -- (opensim);
\draw[arrow] (opensim) -- (clinical);
\draw[arrow] (vicon_ik) -- (validate);
\draw[arrow] (clinical) -- (validate);
\draw[arrow, dashed, darkred] (validate) -- node[right, font=\scriptsize, text=darkred] {정확도 평가} (clinical);
\end{tikzpicture}
\caption{3DGS + Vicon 통합 투구 분석 파이프라인. Vicon 결과는 골드 스탠다드로서 3DGS 기반 결과를 검증하는 데 사용됩니다.}
\label{fig:pitching_pipeline}
\end{figure}

\section{단계별 상세 설명}

\subsection{Step 1: 다시점 2D 자세 추정}

7대 카메라의 각 프레임에서 독립적으로 2D 관절 위치를 추정합니다. 240\,fps의 고속 영상에서는 \textbf{ViTPose} 또는 \textbf{RTMPose}가 OpenPose보다 적합합니다---모션 블러에 더 강건하고, transformer 기반 구조가 전역 문맥을 잘 활용합니다.

\subsection{Step 2: 다시점 SMPL-X 피팅}

\textbf{EasyMocap} 프레임워크로 다시점 2D 키포인트를 종합하여 SMPL-X 파라미터를 추정:

\begin{equation}
(\vtheta^*, \vbeta^*) = \arg\min_{\vtheta, \vbeta} \sum_{j=1}^{7} \sum_{k=1}^{K} w_{jk} \| \mathbf{u}_{jk} - \pi_j(\mathbf{J}_k(\vtheta, \vbeta)) \|^2 + \lambda_\theta E_\theta + \lambda_\beta E_\beta
\label{eq:multiview_smpl}
\end{equation}

여기서:
\begin{itemize}[leftmargin=*]
  \item $j$: 카메라 인덱스 (1--7)
  \item $k$: 관절 인덱스 (1--$K$)
  \item $\mathbf{u}_{jk}$: 카메라 $j$에서 관절 $k$의 2D 검출 좌표
  \item $\pi_j$: 카메라 $j$의 투영 함수
  \item $\mathbf{J}_k(\vtheta, \vbeta)$: SMPL-X 모델에서 관절 $k$의 3D 위치
  \item $w_{jk}$: 가중치 (검출 신뢰도, 가려진 관절은 0)
\end{itemize}

\begin{keyidea}
7대 카메라의 핵심 이점: 단안 영상에서는 깊이 모호성이 심각하지만, 7시점에서는 $\binom{7}{2} = 21$개의 에피폴라 제약이 깊이를 강력하게 결정합니다. 특히 투구 팔의 깊이---앞뒤로 움직이는 정도---를 정확히 복원할 수 있으며, 이는 어깨 외회전/내회전 추정에 결정적입니다.
\end{keyidea}

\subsection{Step 3: 3DGS 인체 재구성 + 자세 정제}

GART 또는 ExAvatar 프레임워크로 캐노니컬 3D 가우시안을 생성하고, 매 프레임 SMPL 자세로 구동합니다.

\textbf{핵심: 렌더링 기반 자세 정제}

초기 SMPL 피팅이 부정확해도, 3DGS의 \textbf{광도 손실(photometric loss)}이 자세를 자동으로 교정합니다:

\begin{equation}
\vtheta^{**} = \arg\min_\vtheta \sum_{j=1}^{7} \Big[ (1-\lambda)\| I_j - \hat{I}_j(\vtheta) \|_1 + \lambda \cdot \mathcal{L}_{\text{D-SSIM}}(I_j, \hat{I}_j(\vtheta)) \Big]
\label{eq:pose_refinement}
\end{equation}

``렌더링된 이미지가 7대 카메라의 실제 영상과 일치하려면 자세가 정확해야 한다''는 강력한 제약이 자세를 정제합니다. 이는 2D 키포인트 재투영만 사용하는 기존 방법보다 \textbf{훨씬 풍부한 감독 신호(supervision)}를 제공합니다---전체 픽셀 수준의 외형 일치를 요구하기 때문입니다.

\subsection{Step 4: 가상 마커 추출 + OpenSim}

정제된 SMPL-X 메시에서 해부학적 랜드마크에 해당하는 꼭짓점을 추출하여 OpenSim IK/ID에 입력합니다. Part 2의 \cref{sec:virtual_markers}와 동일한 파이프라인입니다.

\subsection{Step 5: Vicon 결과와 비교 검증}

Vicon의 마커 기반 IK/ID 결과를 골드 스탠다드로 사용하여 3DGS 기반 결과의 정확도를 평가합니다:

\begin{equation}
\text{RMSE}_{\theta_k} = \sqrt{\frac{1}{T}\sum_{t=1}^{T} (\theta_k^{\text{3DGS}}(t) - \theta_k^{\text{Vicon}}(t))^2}
\end{equation}


\section{3DGS가 투구 분석에 가져오는 고유한 이점}
\label{sec:3dgs_advantages}

\begin{enumerate}[leftmargin=*]
  \item \textbf{마커 가려짐 문제 해결}: 3DGS는 물리적 마커가 아닌 학습된 가우시안으로 표현하므로, 와인드업/팔로스루의 가려짐이 근본적으로 없음

  \item \textbf{연조직 역학 시각화}: 마커는 피부 위의 점일 뿐이지만, 3DGS는 실제 표면 외형을 재구성---삼각근 수축, 전완 회내 시 근육 팽창, 유니폼 변형 등을 시각적으로 관찰 가능

  \item \textbf{임의 시점 렌더링}: 투수의 정면, 후면, 포수 시점, 3루측 시점 등 어떤 각도에서든 고품질 렌더링. 코치와 의료진에게 다각도에서 동작을 보여줄 수 있음

  \item \textbf{사후 분석 가능}: 마커 없이 촬영한 과거 영상도 분석 가능---경기 중 영상 활용 가능성

  \item \textbf{자연스러운 동작}: 마커가 고유수용감각(proprioception)이나 역학에 영향을 줄 수 있다는 우려 없음

  \item \textbf{피로 효과 시각화}: 이닝별 투구 동작의 3D 모델을 나란히 비교하여 피로에 따른 동작 변화를 직관적으로 파악
\end{enumerate}


\section{기술적 도전과 해결 전략}
\label{sec:challenges}

\begin{center}
\renewcommand{\arraystretch}{1.4}
\begin{tabular}{p{4cm}p{8cm}}
\toprule
\textbf{도전} & \textbf{해결 전략} \\
\midrule
30\,ms 가속 단계에서 7 프레임만 확보 & 인접 프레임의 3DGS를 초기값으로 사용(warm-start); 시간적 정규화; Vicon 250/500\,Hz 데이터로 보간 \\
\addlinespace
모션 블러 & 짧은 노출 시간 (1/1000초 이하); 충분한 조명; 이벤트 카메라(EventSplat) 결합 가능성 \\
\addlinespace
손가락/공의 세밀한 움직임 & SMPL-X의 손 관절(30개) 활용; 공에 별도 트래킹 \\
\addlinespace
7대 카메라로 충분한 커버리지? & 투수 주변 $\sim$270$^\circ$ 커버리지로 대부분의 시점에서 3+ 카메라 가시; 가려짐이 적은 배치 설계 \\
\addlinespace
계산 비용 & 프레임별 독립 3DGS는 과도; GART의 캐노니컬 모델 + 프레임별 자세만 최적화로 효율화 \\
\bottomrule
\end{tabular}
\end{center}


% ====================================================================
% CHAPTER 6: 임상 응용
% ====================================================================
\chapter{임상 응용: 부상 예방과 복귀}
\label{ch:clinical}

\section{UCL 부상 위험 예측}
\label{sec:ucl_prediction}

2026년 Bright 등의 획기적 연구: 방송 영상(단안!)만으로 18개 생체역학 지표를 추출하고, 토미 존 수술 예측에 성공했습니다.

\begin{cvbox}[title={방송 영상 기반 부상 예측 (Bright et al., 2026)}]
\begin{itemize}[leftmargin=*]
  \item \textbf{입력}: MLB 방송 영상 (단안, 30\,fps)
  \item \textbf{방법}: DreamPose3D + 생체역학 정제 스택 (뼈 길이 제약, 관절 제한 IK, 스무딩, 대칭성)
  \item \textbf{결과}: 18개 지표 중 16개에서 MAE $< 1^\circ$ (156투 검증)
  \item \textbf{부상 예측}:
  \begin{itemize}
    \item \textbf{토미 존 수술 예측: AUC = 0.811}
    \item \textbf{심각한 팔 부상 예측: AUC = 0.825}
    \item 7,348명 투수 데이터로 검증
  \end{itemize}
\end{itemize}
\end{cvbox}

\begin{importantnote}
방송 영상(30\,fps, 단안)에서도 AUC 0.811이 가능하다면, \textbf{우리의 세팅(240\,fps, 7시점)}에서는 훨씬 높은 정확도가 기대됩니다. 3DGS 기반 자세 정제가 추가되면, 임상적으로 의미 있는 수준의 개인별 부상 위험 프로파일을 생성할 수 있을 것입니다.
\end{importantnote}


\section{피로 모니터링}

투구 수 누적에 따른 생체역학적 변화:

\begin{center}
\renewcommand{\arraystretch}{1.3}
\begin{tabular}{lp{7cm}}
\toprule
\textbf{피로 지표} & \textbf{변화} \\
\midrule
구속 감소 & $\sim$2\,m/s ($\sim$5\,mph) 감소 (말기 이닝) \\
MER 감소 & 어깨 최대 외회전 유의하게 감소 \\
무릎 굴곡 감소 & 더 직립된 착지 $\rightarrow$ 에너지 흡수 능력 저하 \\
팔꿈치 굴곡 감소 & 관절 부하 증가 \\
고관절 굴곡 증가 & 보상적 몸통 자세 변화 \\
구종 변화 & 패스트볼 감소, 변화구 증가 \\
\bottomrule
\end{tabular}
\end{center}

\begin{injurybox}
상습적인 팔 피로 상태에서 투구하는 투수는 부상 위험이 \textbf{36배} 높습니다. 피로 모니터링은 선수의 자각 증상에만 의존할 수 없으며, 객관적 생체역학 측정이 필수적입니다.
\end{injurybox}

\subsection{3DGS 기반 피로 모니터링의 장점}

\begin{enumerate}[leftmargin=*]
  \item 이닝별 동일 투수의 3D 모델을 생성하여 시각적으로 비교
  \item MER, 스트라이드 길이, 몸통 기울기 등의 변화를 자동으로 추적
  \item 코치/트레이너에게 실시간에 가까운 피드백 가능 (학습 시간 $\sim$수 분)
\end{enumerate}


\section{복귀 프로토콜과 재활 모니터링}

\subsection{UCL 재건술 후 복귀}

\begin{center}
\renewcommand{\arraystretch}{1.3}
\begin{tabular}{lcl}
\toprule
\textbf{단계} & \textbf{시기 (개월)} & \textbf{내용} \\
\midrule
초기 재활 & 0--2 & 관절 가동범위 회복, 가벼운 근력 운동 \\
진행성 근력 & 2--4 & 점진적 저항 운동, 고유수용감각 훈련 \\
투구 시작 & $\sim$5 & 짧은 거리 평지 던지기 시작 \\
거리 증가 & 5--8 & 30--60--90--120피트 점진적 증가 \\
마운드 복귀 & 8--10 & 마운드에서 점진적 강도 증가 \\
모의 경기 & 10--12 & 타자 상대 시뮬레이션 \\
경기 복귀 & \textbf{12--14} & 실전 경기 (재래식 UCL 재건) \\
 & \textbf{6--9} & 실전 경기 (내부 보강 수리술) \\
\bottomrule
\end{tabular}
\end{center}

\subsection{3DGS 기반 복귀 모니터링}

재활 과정에서 주기적으로 7대 카메라 세팅에서 투구를 촬영하면:
\begin{itemize}[leftmargin=*]
  \item 부상 전과 동일한 동작 패턴으로 복귀했는지 3D로 비교
  \item 보상 동작(compensation)의 발생 여부를 다각도에서 관찰
  \item 팔꿈치 외반 토크가 안전 범위 내인지 확인
  \item 시계열 3D 모델로 복귀 과정을 문서화
\end{itemize}


\section{웨어러블 센서와의 비교}

\begin{center}
\renewcommand{\arraystretch}{1.3}
\small
\begin{tabular}{p{3cm}cccc}
\toprule
\textbf{시스템} & \textbf{정확도} & \textbf{실시간} & \textbf{경기 중} & \textbf{비용} \\
\midrule
Vicon (마커) & 최고 & $\times$ & $\times$ & \$100K+ \\
KinaTrax & 높음 & $\bigcirc$ & $\bigcirc$ & \$50K+ \\
PULSEthrow (IMU) & $\sim$95\% & $\bigcirc$ & $\bigcirc$ & \$150 \\
본 세팅 (7cam+3DGS) & 높음 (검증 필요) & 준실시간 & 불펜 & \$5--10K \\
방송 영상 (Bright) & 양호 & $\bigcirc$ & $\bigcirc$ & $\sim$\$0 \\
\bottomrule
\end{tabular}
\end{center}


% ====================================================================
% CHAPTER 7: 최신 연구 동향
% ====================================================================
\chapter{최신 연구 동향 (2024--2026)}
\label{ch:trends}

\section{AI/ML 기반 투구 분석}

\begin{enumerate}[leftmargin=*]
  \item \textbf{Bright et al. (2026)}: 방송 영상에서 투구 생체역학 추출 + 토미 존 예측 (AUC 0.811)
  \item \textbf{자동화 투구 단계 분류 (2025)}: 스마트폰 기반 자세 추정 AI로 5단계 투구 분류
  \item \textbf{AIRADL 프레임워크}: 딥러닝 기반 운동선수 부상 위험 평가
\end{enumerate}

\section{다시점 스포츠 재구성}

\begin{enumerate}[leftmargin=*]
  \item \textbf{GS-SFS (IEEE TMM, 2024)}: 3DGS + Shape-from-Silhouette로 대규모 스포츠 장면(농구)의 다인물 재구성
  \item \textbf{M3GYM (CVPR, 2025)}: 8대 카메라, 50+명 피험자, 4,700만 프레임의 다시점 자세 데이터셋
  \item \textbf{SV-GS (2025)}: 스켈레톤 구동 가우시안 스플래팅으로 희소 시점 4D 재구성
  \item \textbf{EventSplat (CVPR, 2025)}: 이벤트 카메라 기반 3DGS---고속 동작의 모션 블러 해결
\end{enumerate}

\section{물리 인식 인체 재구성}

\begin{enumerate}[leftmargin=*]
  \item \textbf{PhysPT (CVPR, 2024)}: 오일러-라그랑주 물리 법칙을 손실 함수에 통합
  \item \textbf{MultiPhys (CVPR, 2024)}: 다인물 물리 인식 3D 동작 추정
  \item \textbf{PhysHMR (2025)}: 휴머노이드 제어 정책 학습으로 물리 기반 메시 복원
  \item \textbf{AddBiomechanics (ECCV, 2024)}: 273명, 70+시간의 모캡+힘 데이터셋---학습 데이터 병목 해소
\end{enumerate}

\section{이 연구들이 우리 세팅에 주는 시사점}

\begin{summary}
\begin{enumerate}[leftmargin=*]
  \item \textbf{방송 영상에서도 부상 예측 가능} $\Rightarrow$ 7시점 240\,Hz에서는 훨씬 더 정확한 예측 가능
  \item \textbf{물리 법칙 통합}이 자세 추정 정확도를 크게 향상 $\Rightarrow$ 3DGS 학습에 물리 손실 추가 가능
  \item \textbf{대규모 쌍 데이터셋} 등장 $\Rightarrow$ 딥러닝 기반 GRF/토크 추정의 학습 데이터 확보
  \item \textbf{이벤트 카메라}가 고속 동작의 모션 블러를 해결 $\Rightarrow$ 향후 세팅에 이벤트 카메라 추가 가능성
  \item \textbf{스포츠 특화 다시점 재구성} 활발 $\Rightarrow$ 야구 투구에 특화된 3DGS 파이프라인 개발의 학문적 기반
\end{enumerate}
\end{summary}


% ====================================================================
% CHAPTER 8: 연구 설계
% ====================================================================
\chapter{연구 설계 제안}
\label{ch:study_design}

\section{연구 질문}

\begin{enumerate}[leftmargin=*]
  \item 7시점 240\,Hz 3DGS 기반 자세 정제가 Vicon 대비 어느 수준의 투구 생체역학 정확도를 달성하는가?
  \item 3DGS 기반 팔꿈치 외반 토크 추정이 임상적으로 유의미한 정확도에 도달하는가?
  \item 3DGS의 포토리얼리스틱 렌더링이 코치/의료진의 동작 판단에 부가 가치를 제공하는가?
\end{enumerate}

\section{실험 프로토콜}

\begin{enumerate}[leftmargin=*]
  \item \textbf{피험자}: 대학 투수 15--20명 (다양한 투구 유형)
  \item \textbf{촬영}: 각 투수 3구종 $\times$ 10투 = 30투/세션
  \item \textbf{장비}: 7대 RGB (240\,Hz) + 10--12대 Vicon (250/500\,Hz) + 힘 측정판 + 공 2개 마커
  \item \textbf{보정}: Vicon 활성 완드 + 체커보드로 RGB-Vicon 좌표계 정합
  \item \textbf{마커}: ISB 권고 상지 마커 세트 (38--59개)
\end{enumerate}

\section{분석 계획}

\begin{center}
\renewcommand{\arraystretch}{1.3}
\begin{tabular}{lp{8cm}}
\toprule
\textbf{비교} & \textbf{지표} \\
\midrule
\textbf{자세 정확도} & 관절 각도 RMSE, 관절 위치 MPJPE (Vicon 기준) \\
\textbf{운동역학 정확도} & 팔꿈치 외반 토크, 어깨 압축력 ($\pm$ 5\Nm 이내?) \\
\textbf{이벤트 검출} & MER, FC, BR 시점 오차 (프레임 수) \\
\textbf{임상 일치도} & Bland-Altman 분석, ICC, SEM \\
\textbf{시각적 평가} & 전문가 블라인드 평가 (5점 리커트) \\
\bottomrule
\end{tabular}
\end{center}

\section{기대 성과}

\begin{itemize}[leftmargin=*]
  \item 7시점 3DGS 기반 투구 분석의 첫 번째 체계적 검증 데이터
  \item Vicon 없이도 임상 수준 투구 분석 가능성 입증 (또는 한계 규명)
  \item 3DGS의 포토리얼리스틱 렌더링 + 생체역학 = 새로운 ``시각적 생체역학'' 보고서 형식 제안
  \item 오픈소스 파이프라인 공개 $\rightarrow$ 후속 연구 촉진
\end{itemize}


% ====================================================================
% Appendix
% ====================================================================
\chapter{수학적 부록}
\label{ch:appendix}

\section{역동역학으로 팔꿈치 외반 토크 계산}
\label{sec:valgus_torque_calc}

전완 분절에 대한 뉴턴-오일러 방정식 (발단 $\rightarrow$ 근위):

\begin{align}
\vF_{\text{elbow}} &= m_{\text{forearm}} \va_{\text{CoM}} - m_{\text{forearm}} \vg - \vF_{\text{wrist}} - \vF_{\text{ball}} \\
\vM_{\text{elbow}} &= \mI_{\text{forearm}} \bm{\alpha} + \bm{\omega} \times (\mI_{\text{forearm}} \bm{\omega}) - \vM_{\text{wrist}} \notag \\
&\quad + \vr_{\text{CoM} \rightarrow \text{wrist}} \times \vF_{\text{wrist}} - \vr_{\text{CoM} \rightarrow \text{elbow}} \times \vF_{\text{elbow}}
\end{align}

$\vM_{\text{elbow}}$의 내외측 성분이 \textbf{외반/내반 토크}입니다. 이 값이 양수(내반, varus)이면 UCL을 보호하고, 음수(외반, valgus)이면 UCL에 부하를 가합니다.

\textbf{입력 데이터 요구사항}:
\begin{itemize}[leftmargin=*]
  \item 전완 CoM 가속도 $\va_{\text{CoM}}$: 위치의 이중 미분 (SMPL 또는 마커로부터)
  \item 각속도 $\bm{\omega}$, 각가속도 $\bm{\alpha}$: 회전행렬의 시간 미분
  \item 손목 관절력 $\vF_{\text{wrist}}$: 손 분절의 역동역학에서 계산
  \item 공에 의한 힘 $\vF_{\text{ball}}$: 릴리스 전에만 작용 (공 가속도 $\times$ 공 질량)
\end{itemize}

\begin{importantnote}
팔꿈치 외반 토크 계산에서 가장 민감한 변수는 \textbf{관절 중심 위치}입니다. 관절 중심이 5\,mm 이동하면 모멘트 팔이 변하여 토크 추정치가 수 Nm 변합니다. 이것이 마커리스 시스템의 관절 위치 오차(52--57\,mm)가 투구 분석에서 심각한 문제인 이유입니다.
\end{importantnote}


\section{다시점 삼각측량 정밀도}
\label{sec:triangulation_precision}

$N$대 카메라에서의 삼각측량 정밀도. 각 카메라 $j$에서의 2D 검출 오차가 $\sigma_{\text{2D}}$일 때, 3D 복원 오차:

\begin{equation}
\sigma_{\text{3D}} \approx \frac{\sigma_{\text{2D}} \cdot d}{f \cdot B \cdot \sqrt{N-1}}
\end{equation}

여기서:
\begin{itemize}[leftmargin=*]
  \item $d$: 피사체까지의 거리
  \item $f$: 초점 거리 (픽셀 단위)
  \item $B$: 카메라 간 평균 기선 길이(baseline)
  \item $N$: 카메라 수
\end{itemize}

7대 카메라($N=7$)는 2대($N=2$) 대비 $\sqrt{6}/\sqrt{1} \approx 2.45$배 정밀합니다.


\section{SMPL 관절에서 해부학적 관절 각도 변환}
\label{sec:smpl_to_clinical}

SMPL의 축-각도 $\vtheta_k \in \RR^3$를 회전행렬 $\mR_k$로 변환 후, ISB 권고 분해 순서에 따라:

\textbf{견관절 (Y-X-Y)}:
\begin{align}
\text{거상 평면} &= \text{atan2}(R_{31}, R_{11}) \\
\text{거상각} &= \text{acos}(R_{21}) \\
\text{축 회전} &= \text{atan2}(R_{23}, -R_{22})
\end{align}

\textbf{주관절 (Z-X-Y)}:
\begin{align}
\text{굴곡/신전} &= \text{atan2}(-R_{23}, R_{33}) \\
\text{운반각} &= \text{asin}(R_{13}) \\
\text{회내/회외} &= \text{atan2}(-R_{12}, R_{11})
\end{align}

\begin{remark}
SMPL은 모든 관절을 3-DOF로 모델링하므로, 주관절의 ``운반각'' 성분이 비현실적으로 크게 나올 수 있습니다. 생체역학적 관절 제한(joint constraint)을 적용하여 이를 억제해야 합니다.
\end{remark}


% ==================== REFERENCES ====================
\chapter*{참고 문헌}
\addcontentsline{toc}{chapter}{참고 문헌}

\begin{enumerate}[leftmargin=*, label={[\arabic*]}]
  \item Fleisig, G.S. et al. ``Kinetics of Baseball Pitching with Implications About Injury Mechanisms.'' \textit{Am. J. Sports Med.}, 23(2), 1995.

  \item Diffendaffer, A.Z. et al. ``The Clinician's Guide to Baseball Pitching Biomechanics.'' \textit{Sports Health}, 15(2), 2023.

  \item Escamilla, R.F. et al. ``Pitching Biomechanics as a Pitcher Approaches Muscular Fatigue During a Simulated Baseball Game.'' \textit{Am. J. Sports Med.}, 35(1), 2007.

  \item Wu, G. et al. ``ISB Recommendation on Definitions of Joint Coordinate Systems for the Shoulder, Elbow, Wrist and Hand.'' \textit{J. Biomech.}, 38(5), 2005.

  \item Bright, J., Mende, K. \& Zelek, J. ``Scalable Injury-Risk Screening in Baseball Pitching From Broadcast Video.'' arXiv:2603.04864, 2026.

  \item KinaTrax Validation Study. ``Normative In-Game Pitching Biomechanics Data Using a Markerless Motion-Capture System.'' \textit{Am. J. Sports Med.}, 2023.

  \item Assessment of Markerless Accuracy for Baseball Pitching. \textit{J. Sports Sciences}, 2025.

  \item Lei, J. et al. ``GART: Gaussian Articulated Template Models.'' \textit{CVPR}, 2024.

  \item Moon, G. et al. ``ExAvatar: Expressive Whole-Body 3D Gaussian Avatar.'' \textit{ECCV}, 2024.

  \item Kerbl, B. et al. ``3D Gaussian Splatting for Real-Time Radiance Field Rendering.'' \textit{ACM TOG (SIGGRAPH)}, 2023.

  \item Wu, G. et al. ``4D Gaussian Splatting for Real-Time Dynamic Scene Rendering.'' \textit{CVPR}, 2024.

  \item Yura, A. et al. ``EventSplat: 3D Gaussian Splatting from Moving Event Cameras.'' \textit{CVPR}, 2025.

  \item Zhang, J. et al. ``PhysPT: Physics-aware Pretrained Transformer.'' \textit{CVPR}, 2024.

  \item Werling, K. et al. ``AddBiomechanics Dataset.'' \textit{ECCV}, 2024.

  \item Loper, M. et al. ``SMPL: A Skinned Multi-Person Linear Model.'' \textit{ACM TOG}, 2015.

  \item Dillman, C.J. et al. ``Biomechanical Aspects of Pitching.'' \textit{Am. J. Sports Med.}, 21(2), 1993.
\end{enumerate}


% ==================== GLOSSARY ====================
\chapter*{용어집}
\addcontentsline{toc}{chapter}{용어집}

\begin{description}[leftmargin=!, labelwidth=5cm, font=\normalfont\bfseries]
  \item[Arm Cocking] 팔 코킹. 투구 팔이 외회전하며 에너지를 저장하는 단계.
  \item[Ball Release (BR)] 공 릴리스. 공이 손에서 떠나는 순간.
  \item[Baseline] 기선. 두 카메라 사이의 거리. 기선이 길수록 삼각측량 정밀도 향상.
  \item[Extrinsic Parameters] 외부 파라미터. 카메라의 3D 공간 위치와 방향.
  \item[Flexor-Pronator Mass] 굴근-회내근 복합체 (FPM). 팔꿈치 내측의 근육군으로 UCL을 보호.
  \item[Foot Contact (FC)] 앞발 접지. 투수의 앞발이 지면에 닿는 순간.
  \item[Genlock] 동기 잠금. 모든 카메라를 동일 클럭 신호에 동기화하는 기술.
  \item[Hip-Shoulder Separation] 골반-몸통 분리. 앞발 접지 시 골반과 어깨의 회전각 차이.
  \item[Intrinsic Parameters] 내부 파라미터. 카메라의 초점거리, 주점, 왜곡 계수.
  \item[KinaTrax] MLB에서 사용하는 경기장 기반 마커리스 모션캡처 시스템.
  \item[MER] Maximum External Rotation. 최대 외회전. 투구 팔이 뒤로 가장 많이 회전한 시점.
  \item[Moment Arm] 모멘트 팔. 힘의 작용선과 관절 중심 사이의 수직 거리.
  \item[MPJPE] Mean Per Joint Position Error. 관절별 평균 위치 오차 (mm).
  \item[Nyquist Theorem] 나이퀴스트 정리. 신호를 정확히 복원하려면 최소 2배의 샘플링 주파수 필요.
  \item[Photometric Loss] 광도 손실. 렌더링된 이미지와 실제 이미지의 픽셀 수준 차이.
  \item[Pitch Count] 투구 수. 부상 예방을 위한 투구 횟수 제한 관리.
  \item[Pose Refinement] 자세 정제. 렌더링 기반 최적화로 초기 자세 추정을 교정하는 과정.
  \item[Spin Efficiency] 회전 효율. 총 회전 중 공 궤적에 유효한 성분의 비율.
  \item[Stride Length] 스트라이드 길이. 투구판에서 앞발 착지 지점까지의 거리.
  \item[Tommy John Surgery] 토미 존 수술. UCL 재건술의 통칭.
  \item[UCL] Ulnar Collateral Ligament. 척골측부인대. 팔꿈치 내측의 주요 인대.
  \item[Valgus Torque] 외반 토크. 팔꿈치를 바깥으로 벌어지게 하는 토크. UCL 부하의 핵심 지표.
  \item[Varus Torque] 내반 토크. 외반 토크에 대항하는 보호 토크 (근육에 의해 생성).
  \item[Warm-Starting] 이전 프레임의 3DGS 결과를 현재 프레임의 초기값으로 사용하는 전략.
\end{description}


\end{document}
